\begin{abstract}
\noindent This note documents the TOFME project, whose goal is a mixed-event
estimate of random (fake) matches between TPC tracks and Barrel Time-of-Flight (BTOF) hits in STAR.
The input is the latest Run19 Au+Au 19.6 GeV production (P24iy, library SL24y).
First, a QA analysis of the picoDst content was set up and run:
event-level, TPC-track-level and TPC--BTOF-association variables, including $1/\beta$ and
\mtwo{} versus momentum, local hit coordinates, tray/module occupancy and the TOF matching efficiency.
Then the mixed-event machinery was built (Sec.~\ref{sec:me}): a data-driven BTOF module geometry, projection of
the picoDst helix of every track to a BTOF cell, and cell-level matching to hits of the same and of mixed events.
For the production-like matching rule (projected cell or its neighbour), a random hit is present in the matching cells
for \mePMENbpctVthree\% of the tracks, equal to the hit occupancy (1M-event sample, \meNRunsVthree{} runs,
\meEvtSelVthree{} selected events; stable from run to run once scaled by the occupancy; Sec.~\ref{sec:me1m}).
Finally (Sec.~\ref{sec:m2}), $\beta$ and \mtwo{} are computed for same-event and mixed-event matches in nine
centrality classes. The mixed events are overlaid on the track's own event, normalized per projected track, and include
hits without a stored time. Random matches are then \mtMENormFracCiVthree\% of the TOF-matched tracks in 0--5\% and
\mtMENormFracCaVthree\% in 70--80\%, and \mtFrPiMidCiVthree\% ($\pi$), \mtFrKMidCiVthree\% (K) and \mtFrPMidCiVthree\% (p) of
the yield under the PID peaks in 0--5\% at $0.8<p<1$\,\GeVc. They explain only a small part of the yield between the
proton and deuteron peaks, because the picoDst has no time for BTOF hits not matched to a track.
On MuDst (Sec.~\ref{sec:mudst}), where all hits have a time, the production BTOF matching and calibration re-run on MuDst
reproduce production exactly. They are then re-run on overlaid events (hits and tracks of a partner event added); this
mixed event describes the same-event background between the p and d peaks (SE/ME = \ovtHiCi{} in 0--5\%, against
\mtSbRatioAllVnine{} on picoDst). Random matches are \ovtFrCi\% of the TOF matches in 0--5\% and \ovtFrAll\% in 0--80\%, in
agreement with an independent emulation (\muFrProdCi\% and \muFrProdAll\%; four MuDst files of two runs).
A HIJING+GEANT simulation of the 2019 detector (Sec.~\ref{sec:sim}, \simEvGen{} events) shows that the remaining
low-momentum excess is track-correlated: at 0.2--0.4\,GeV/$c$ it is mostly secondaries produced by the track's own particle in the
BTOF trays, at a rate (\simCorrA{} per $10^4$ tracks) flat in centrality as in the data. The simulated component over-predicts the K--p valley
and is not used as a correction.
\textbf{Limitations:} the QA sections (\ref{sec:evtqa}--\ref{sec:tofqa}) use a \emph{test sample} of \nProcEvents{} events from one
picoDst file of a single run (20093036); the mixed-event results use 15 runs. No bad-tray list is applied.
The time of random hits without a stored time is an assumption in the picoDst sections (Sec.~\ref{sec:m2}); the MuDst
results use \ovEvMtwo{} events only. Statistical uncertainties only.
Coloured boxes: \textcolor{blue!60!black}{\textbf{blue}} = a choice made and its reason;
\textcolor{red!70!black}{\textbf{red}} = caveat / not yet understood;
\textcolor{orange!80!black}{\textbf{orange}} = decision needed from the analyser.
\end{abstract}
